\section{Deskripsi Persoalan}
\subsection{Latar Belakang dan Tujuan}
Pemuatan paket pada truk merupakan persoalan optimasi yang melibatkan pemilihan
paket dan pengaturan letaknya dalam ruang terbatas. Paket dapat memiliki dimensi,
nilai, berat, serta sifat kerapuhan yang berbeda. Pemilihan berdasarkan nilai
semata belum tentu menghasilkan susunan yang dapat dimuat, karena penempatan
paket juga harus memenuhi batas ruang, kapasitas berat, dan aturan penopangan.

Studi kasus Mejiro Express pada tugas besar ini meminta konfigurasi pemuatan
paket yang memaksimalkan total nilai muatan sebuah truk~\cite{spesifikasi}.
Tidak semua paket harus diangkut. Oleh karena itu, solusi mencakup keputusan
paket mana yang dimuat, posisi paket tersebut, dan orientasinya.

Tujuan pengerjaan adalah merancang representasi solusi, menerapkan salah satu
varian \textit{hill-climbing}, \textit{simulated annealing} (SA), dan
\textit{genetic algorithm} (GA), serta mengevaluasi kualitas solusi, durasi
pencarian, dan konsistensi hasilnya. Evaluasi juga mencakup pengaruh jumlah
populasi dan banyak iterasi pada GA.

\subsection{Definisi Permasalahan dan Data Masukan}
Masukan terdiri atas himpunan paket $P=\{1,\ldots,N\}$ dan sebuah truk dengan
satu peti kemas. Setiap paket berbentuk balok atau kubus. Atribut masukan
dirangkum pada Tabel~\ref{tab:atribut}. Satuan dimensi harus konsisten antarentitas.

\begin{table}[H]\centering\small
\caption{Atribut paket dan truk}\label{tab:atribut}
\begin{tabularx}{\linewidth}{@{}llY@{}}\toprule
Entitas & Atribut & Makna\\\midrule
Paket & \texttt{id} & Identitas paket.\\
 & \texttt{Dimensions} & Lebar, panjang, dan tinggi dasar paket.\\
 & \texttt{Orientation} & Pemetaan sisi paket terhadap sumbu koordinat.\\
 & \texttt{Value} & Nilai paket dalam satuan unit.\\
 & \texttt{Weight} & Berat paket dalam satuan unit berat.\\
 & \texttt{IsFragile} & Penanda paket pecah-belah.\\
 & \texttt{ETA} & Perkiraan waktu kedatangan dalam sisa hari.\\\midrule
Truk & \texttt{Dimensions} & Lebar, panjang, dan tinggi peti kemas.\\
 & \texttt{MaxCapacity} & Kapasitas maksimum total berat muatan.\\\bottomrule
\end{tabularx}\end{table}

Ukuran problem set, jumlah paket, distribusi dimensi, berat, dan komposisi paket
dibebaskan kepada kelompok selama sesuai spesifikasi~\cite{qna}. Kebebasan ini
disertai kebutuhan mendokumentasikan format input dan syarat validitasnya.
Tidak ada format test case tetap yang diumumkan asisten pada QnA yang tersedia.
Pernyataan tersebut tidak menjamin ketiadaan test case tambahan saat demo.

\subsection{Ketentuan dan Batasan Persoalan}
Ketentuan pemodelan yang menjadi acuan adalah sebagai berikut~\cite{spesifikasi,qna}.
\begin{enumerate}
\item Ruang lingkup wajib mencakup satu truk dengan satu peti kemas.
\item Paket dapat berbeda ukuran, dan tidak seluruh paket harus dimuat.
      Paket di luar harus seluruhnya berada di luar; paket tidak boleh memotong dinding.
\item Posisi paket menggunakan koordinat bilangan bulat. Initial state dibentuk acak.
\item Orientasi menyatakan pemetaan sisi paket ke sumbu; tersedia enam kemungkinan
      orientasi. Langkah rotasi neighbor adalah $90^\circ$ terhadap sumbu $x$, $y$, atau $z$.
\item Paket di dalam tidak boleh melayang. Adanya support pada sebagian alas
      sudah cukup menurut definisi spesifikasi dan klarifikasi asisten.
\item Paket fragile tidak boleh menjadi support paket lain.
\item Jumlah berat paket di dalam tidak boleh melebihi kapasitas truk.
\item Neighbor dibentuk melalui pertukaran, pemindahan, atau rotasi paket.
      Pemindahan dapat mengubah posisi paket dari dalam ke luar truk maupun sebaliknya.
\end{enumerate}

\subsection{Tujuan Optimasi dan Keluaran Program}
Objektif wajib adalah memaksimalkan jumlah nilai paket yang berada di dalam
peti kemas. Solusi akhir harus memperlihatkan status semua paket, termasuk
paket yang tidak dimuat, posisi dan orientasi paket yang dimuat, keadaan ruang
peti kemas, dan nilai objektif state. Visualisasi program juga harus dapat
menampilkan initial state, final state, dan hasil eksperimen.

\subsection{Klarifikasi dan Asumsi Pemodelan}
Pertukaran ukuran efektif akibat rotasi merupakan perubahan orientasi, bukan
perubahan dimensi fisik paket. Sebagai contoh, dimensi dasar $(3,4,1)$ dapat
memiliki dimensi efektif $(4,3,1)$ setelah rotasi terhadap sumbu $z$. Apabila
dua atau tiga dimensi sama, orientasi berbeda dapat menghasilkan ukuran
efektif yang sama; pencatatan orientasi perlu tetap konsisten.

QnA memperbolehkan kelompok menentukan kegunaan ETA dengan justifikasi.
Pada rancangan solusi ini, ETA digunakan sebagai informasi urgensi untuk
heuristik dan/atau objektif bonus. ETA yang lebih kecil merepresentasikan
waktu kedatangan yang lebih dekat. Pemanfaatannya dibedakan dari fungsi
objektif wajib yang tetap memaksimalkan total nilai paket di dalam truk.
Aturan prioritas berbasis ETA tidak menghapus kewajiban memenuhi seluruh
constraint geometris, fragile, dan kapasitas berat.


Larangan tumpang tindih volume paket diperlakukan sebagai asumsi geometris
untuk menghasilkan susunan fisik yang masuk akal. Spesifikasi tidak
mencantumkannya sebagai butir constraint tersendiri. Asumsi ini perlu
dipertahankan secara konsisten dalam validasi dan dijelaskan saat demo.
Tidak ditambahkan kewajiban support seluruh alas, pemeriksaan titik berat,
ataupun batas beban tumpukan karena ketentuan tersebut tidak diberikan.
